% Part: first-order-logic
% Chapter: introduction
% Section: satisfaction

\documentclass[../../../include/open-logic-section]{subfiles}

\begin{document}

\olfileid[fr]{fol}{int}{sat}

\olsection{Satisfaction}

Dès que nous savons ce que sont les !!{formula}s, nous pouvons déjà
passer à la sémantique de la logique du premier ordre. La notion
fondamentale est alors celle d'!!a{structure}. Pour notre langage
simple, !!a{structure}~$\Struct M$ ne possède que trois composantes :
un ensemble non vide $\Domain{M}$ appelé le \emph{!!{domain}}, l'objet
que $\Obj a$ désigne dans~$\Struct M$, et l'ensemble des objets auxquels
$\Obj P$ s'applique dans~$\Struct M$. L'objet désigné par~$\Obj a$ est
noté~$\Assign{\Obj a}{M}$, et l'ensemble des objets auxquels $\Obj P$
s'applique,~$\Assign{\Obj P}{M}$. !!^a{structure}~$\Struct{M}$ est donc
constituée exactement de ces trois éléments :
$\Domain{M}$, $\Assign{\Obj a}{M} \in \Domain{M}$ et
$\Assign{\Obj P}{M} \subseteq \Domain{M}$. Le cas général sera plus
complexe, puisqu'il y aura de nombreux !!{predicate}s et de nombreuses
!!{constant}s, que les !!{predicate}s pourront avoir plusieurs places
et que le langage comportera aussi des !!{function}s.

Cela suffit pour définir la satisfaction des !!{formula}s qui ne
contiennent pas de !!{variable}. L'idée consiste à donner une définition
inductive qui reproduit la manière dont nous avons défini les
!!{formula}s. Nous indiquons quand une formule atomique est satisfaite
dans~$\Struct{M}$, puis, par exemple, quand $\lnot !A$ est satisfaite
dans~$\Struct{M}$ en fonction de la satisfaction de~$!A$
dans~$\Struct{M}$. Nous
pourrions ainsi poser :
\begin{enumerate}
  \item $\Atom{\Obj P}{\Obj a}$ est satisfaite dans~$\Struct{M}$ si et
  seulement si $\Assign{\Obj a}{M} \in \Assign{\Obj P}{M}$.
  \item $\lnot !A$ est satisfaite dans~$\Struct{M}$ si et seulement si
  $!A$ n'est pas satisfaite dans~$\Struct{M}$.
  \item $(!A \land !B)$ est satisfaite dans~$\Struct{M}$ si et
  seulement si $!A$ est satisfaite dans~$\Struct{M}$ et si $!B$ est
  elle aussi satisfaite dans~$\Struct{M}$.
\end{enumerate}
Posons $\Domain{M} = \{0, 1, 2\}$,
$\Assign{\Obj a}{M} = 1$ et
$\Assign{\Obj P}{M} = \{1, 2\}$. Cette définition nous apprend que
$\Atom{\Obj P}{\Obj a}$ est satisfaite dans~$\Struct{M}$, puisque
$\Assign{\Obj a}{M} = 1 \in \{1,2\} = \Assign{\Obj P}{M}$. Elle nous
apprend ensuite que $\lnot \Atom{\Obj P}{\Obj a}$ n'est pas satisfaite
dans~$\Struct{M}$, que
$\lnot\lnot \Atom{\Obj P}{\Obj a}$ l'est, et que
$(\lnot \Atom{\Obj P}{\Obj a} \land \Atom{\Obj P}{\Obj a})$ ne l'est
pas, et ainsi de suite.

La difficulté apparaît lorsqu'il faut définir le cas des
quantificateurs. Nous voudrions écrire quelque chose comme :
« $\lexists[\Obj v_0][\Atom{\Obj P}{\Obj v_0}]$ est satisfaite si et
seulement si $\Atom{\Obj P}{\Obj v_0}$ est satisfaite. » Mais la
!!{structure}~$\Struct{M}$ ne nous dit pas comment traiter les
!!{variable}s. Ce que nous voulons réellement dire, c'est que
$\Atom{\Obj P}{\Obj v_0}$ est satisfaite \emph{pour une certaine valeur
de~$\Obj v_0$}. Pour rendre cette idée précise, il nous faut un moyen
d'associer des !!{element}s de~$\Domain{M}$ non seulement à $\Obj a$,
mais aussi à~$\Obj v_0$. Nous introduisons à cette fin les
\emph{assignations} de !!{variable}s. Une assignation de
!!{variable}s est simplement une fonction~$s$ qui associe à chaque
!!{variable} un !!{element} de~$\Domain{M}$, donc, dans notre exemple,
l'un des
éléments $0$, $1$ ou~$2$. Comme nous ne savons pas à l'avance quelles
!!{variable}s pourront figurer dans !!a{formula}, nous ne pouvons pas
restreindre les !!{variable}s auxquelles $s$ attribue une valeur. La
solution la plus simple consiste à exiger que $s$ attribue une valeur à
\emph{toutes} les !!{variable}s $\Obj v_0$, $\Obj v_1$, \dots\@; nous
n'utiliserons que celles dont nous avons besoin.

Au lieu de définir la satisfaction des !!{formula}s relativement à une
seule !!{structure}, nous la définirons relativement à
!!a{structure}~$\Struct{M}$ \emph{et} à une assignation de
!!{variable}s~$s$, et nous écrirons en abrégé $\Sat{M}{!A}[s]$. Notre
définition comportera désormais une clause supplémentaire pour les
!!{formula}s atomiques contenant des !!{variable}s :
\begin{enumerate}
  \item $\Sat{M}{\Atom{\Obj P}{\Obj a}}[s]$ si et seulement si
  $\Assign{\Obj a}{M} \in \Assign{\Obj P}{M}$.
  \item $\Sat{M}{\Atom{\Obj P}{\Obj v_i}}[s]$ si et seulement si
  $s(\Obj v_i) \in \Assign{\Obj P}{M}$.
  \item $\Sat{M}{\lnot !A}[s]$ si et seulement si
  $\Sat{M}{!A}[s]$ ne vaut pas.
  \item $\Sat{M}{(!A \land !B)}[s]$ si et seulement si
  $\Sat{M}{!A}[s]$ et $\Sat{M}{!B}[s]$.
\end{enumerate}
Une première difficulté est ainsi résolue : nous pouvons dire quand
$\Struct{M}$ satisfait $\Atom{\Obj P}{\Obj v_0}$ pour la
valeur~$s(\Obj v_0)$. Pour définir correctement la satisfaction de
$\lexists[\Obj v_0][\Atom{\Obj P}{\Obj v_0}]$, une étape supplémentaire
est nécessaire. Nous voulons que
$\Sat{M}{\lexists[\Obj v_0][\Atom{\Obj P}{\Obj v_0}]}[s]$ vaille si et
seulement si
$\Sat{M}{\Atom{\Obj P}{\Obj v_0}}[s']$ vaut pour \emph{une certaine}
assignation~$s'$ qui attribue une valeur à~$\Obj v_0$. Cette valeur
attribuée à~$\Obj v_0$ ne doit toutefois pas nécessairement
être~$s(\Obj v_0)$. Introduisons la
notation suivante : si $m \in \Domain{M}$,
$\Subst{s}{m}{\Obj v_0}$ désigne l'assignation qui coïncide avec~$s$
sur toutes les !!{variable}s autres que~$\Obj v_0$, et qui associe
$m$ à~$\Obj v_0$. Nous pouvons alors ajouter la clause :
\begin{enumerate}\setcounter{enumi}{4}
  \item $\Sat{M}{\lexists[\Obj v_i][!A]}[s]$ si et seulement si
  $\Sat{M}{!A}[\Subst{s}{m}{\Obj v_i}]$ pour un certain
  $m \in \Domain{M}$.
\end{enumerate}
Cette définition donne-t-elle le résultat attendu ? Supposons que
$s(\Obj v_i) = 0$ pour tout~$i \in \Nat$.
$\Sat{M}{\lexists[\Obj v_0][\Atom{\Obj P}{\Obj v_0}]}[s]$ vaut si et
seulement s'il existe $m \in \Domain{M}$ tel que
$\Sat{M}{\Atom{\Obj P}{\Obj v_0}}[\Subst{s}{m}{\Obj v_0}]$.
Un tel $m$ existe : nous pouvons choisir $m = 1$ ou $m = 2$. Remarquons
que l'énoncé est vrai même si la valeur~$s(\Obj v_0)$ que $s$
elle-même attribue à~$\Obj v_0$, ici $0$, ne convient pas. Nous avons
$\Sat{M}{\Atom{\Obj P}{\Obj v_0}}[\Subst{s}{1}{\Obj v_0}]$, mais non
$\Sat{M}{\Atom{\Obj P}{\Obj v_0}}[s]$.

Si tout cela paraît déroutant et lourd, c'est parce que cela l'est. La
complexité supplémentaire est toutefois nécessaire pour donner une
définition inductive précise de la satisfaction de toutes les
!!{formula}s, et il nous faut une définition de ce genre pour rendre
précises les notions sémantiques. D'autres méthodes existent, mais elles
sont toutes aussi peu élégantes.

\begin{editorial}
Deux lapsus de la source sont corrigés. Dans la description du cas
général, ce sont les prédicats, et non les constantes, qui peuvent
avoir plusieurs places. Plus loin, l'assignation prend ses valeurs dans
le domaine $\{0,1,2\}$ déjà fixé, et non dans l'ensemble
$\{1,2,3\}$ indiqué par erreur.
\end{editorial}

\end{document}
